\phantomsection

\subsection*{Sensação de Boca}
\addcontentsline{toc}{subsection}{Sensação de Boca}

\begin{itemize}

\item Os componentes de acidez, taninos, carbonatação, álcool, corpo e dulçor se combinam para formar a \textbf{estrutura} do hidromel, que também pode ser descrita pelo termo cervejeiro \textbf{sensação de boca}.
\item De modo geral, o \textbf{corpo} e a sensação de preenchimento do hidromel aumentam com o dulçor e a força alcoólica. A acidez e a carbonatação frequentemente podem cortar ou reduzir a percepção de corpo. Exemplos extremos (excessivamente leves ou enjoativos) são indesejáveis.
\item Como mencionado em Aparência, o hidromel pode apresentar níveis de \textbf{carbonatação} Tranquilo, Frisante ou Espumante. Níveis mais elevados de carbonatação tendem a tornar a percepção da bebida mais viva, acrescentando acidez, tornando a sensação no paladar mais leve e proporcionando maior frescor e definição ao final.
\item Os \textbf{taninos}, incluindo aqueles provenientes do uso de especiarias, frutas ou de madeira, podem afetar a sensação de boca, aumentando o corpo, o amargor ou a percepção de secura no final, ao acrescentarem uma sensação de repuxamento e persistência na boca. Os taninos também podem acrescentar adstringência e amargor, que podem equilibrar o dulçor.
\item a força alcoólica afeta a sensação de boca principalmente por meio da sensação de aquecimento. O hidromel pode ser classificado de acordo com sua força alcoólica (ABV - alcohol-by-volume) como:

\begin{itemize}

\item \textbf{Leve (Hydromel)}: ABV 3.5 - 7.5\%. Geralmente sem percepção de aquecimento alcoólico. Hydromel (ou suas variações) significa hidromel em diversas línguas românicas. Session é uma denominação alternativa razoável, derivada do vocabulário cervejeiro.
\item \textbf{Moderado (Standard)}: ABV 7.5 - 14.0\%. Semelhante a um vinho de mesa, podendo apresentar leve aquecimento alcoólico, mas sem sensação de queimação.
\item \textbf{Forte (Sack)}: ABV 14.0 - 18.0\%. Geralmente apresenta uma presença alcoólica perceptível, mas não deve ser quente, semelhante a solvente, pungente ou desagradável.

\end{itemize}

\item Níveis mais elevados de dulçor podem mascarar a força alcoólica, portanto hidroméis fortes e doces devem ser avaliados com atenção.

\end{itemize}